% vim: set spell spelllang=es syntax=tex :

\section{Visión por computadora}

\label{mt_visionComputadora}

La disciplina de Visión por Computadora engloba el estudio de métodos de
captura, mejoramiento, y extracción de información a partir de imágenes. Su
objetivo principal es definir un modelo de una escena a partir de
imágenes \cite{cvLinda2001}. En el caso particular del fútbol de robots, el
modelo que se quiere definir describe la posición y orientación de los robots y
la pelota dentro de la cancha.

En computación las imágenes suelen representarse como matrices de vectores de
$N$ componentes. Para las imágenes de escala de grises solo se utiliza un
componente, mientras que para las imágenes a color se utilizan 3 o 4,
dependiendo del espacio de colores. Los espacios de colores mas comunes son
\emph{RGB}, donde los componentes indican las cantidades de rojo, verde y
azul, \emph{YUV}, donde el componente \emph{Y} indica la luma y los
componentes \emph{U} y \emph{V} indican la cromancia, y \emph{HSV} donde sus
componentes indican la tonalidad, saturación, y valor, respectivamente.

Dependiendo de la aplicación, existen distintas etapas para la obtención de la
información y su procesamiento (Fig. \ref{etapasVision}). En general, se definen los
siguientes pasos \cite{digitalImageProcessing2ed,wikiCV}:

\begin{description}

	\item[Adquisición de la imagen:] Se trata de la digitalización de la
		imagen a procesar. Ésta puede ser obtenida de uno o más
		sensores. Distintos sensores pueden capturar distintas bandas
		del espectro electromagnético, e incluso puede tratarse de
		imágenes obtenidas de otros tipos de señales desde sonido
		hasta ondas gravitatorias. El resultado puede variar desde una
		imagen bidimensional monocromática hasta volúmenes
		tridimensionales.

		En el fútbol de robots de tamaño pequeño, el conjunto de cámaras
		captura una serie de imágenes bidimensionales a color. El
		sistema de coordenadas de color de la imagen es normalmente
		\emph{RGB} o \emph{YUV}.

	\item[Preprocesamiento:] Para que la imagen pueda ser procesada, muchas
		veces es necesario aplicarle transformaciones que no extraen
		información, pero que son precondiciones para las siguientes
		etapas de procesamiento. Algunas de las transformaciones comunes
		son corrección de color, reducción de ruido o cambio de espacio
		de color.

		Los sistemas de visión global para el fútbol de robots de la
		\emph{SSL} suelen trabajar con el espacio de color \emph{YUV}
		o \emph{HSV}, por lo que, de ser necesario, la imagen debe ser
		convertida a uno de estos espacios de color. Otro
		preprocesamiento que puede realizarse es la composición de las
		imágenes del conjunto de cámaras en una sola imagen.

	\item[Extracción de características:] Consiste en el reconocimiento de
		características de bajo nivel, como puntos de interés, bordes,
		regiones y masas\footnote{En la bibliografía en ingles se
		utiliza el termino \emph{blob}, que en su uso coloquial
		describe algo carente de forma. En visión por computadora se
		lo suele entender como una región continua de un mismo
		color.}. Los sistemas para la \emph{SSL} suelen en esta etapa
		reconocer las masas de colores de los parches y la pelota.

	\item[Detección y Segmentación:] Se trata de separar la imagen en sus
		partes constituyentes y seleccionar aquéllas de interés. Los
		sistemas de visión global para fútbol de robots de la \emph{SSL}
		utilizan las masas de color reconocidas en la etapa anterior
		para encontrar los robots y la pelota.

	\item[Procesamiento de alto nivel:] En esta etapa se crea finalmente el
		modelo de la escena utilizando los resultados de las etapas
		anteriores. En el fútbol de robots de la \emph{SSL}, se toman
		las posiciones, dentro de la imagen, de los robots y pelota
		encontrados en la etapa anterior, y se calcula su posición
		dentro de la cancha de acuerdo con la geometría de ésta.

	\item[Toma de decisiones:] Muchas de las aplicaciones de visión por
		computadora, en lugar de sólo entregar información, integran la
		toma de decisiones como parte del procesamiento. En el caso de
		los sistemas de visión global para fútbol de robots, el sistema
		sólo notifica a las computadoras de los equipos la posición de
		los robots y de la pelota.

\end{description}

Esta separación de las etapas es principalmente conceptual y distintos sistemas
pueden presentar otra separación o incluso acoplar todas las etapas de
procesamiento en una sola (la etapa de \textbf{adquisición de la imagen} es la
única que no puede ser acoplada con el resto).

\begin{figure}[!ht]

	\centering
	\includegraphics[width=\textwidth]{img/etapasVision.pdf}
	\caption{Etapas de procesamiento de un sistema de visión por
	computadora.}

	\label{etapasVision}

 \end{figure}
